\section{Integer Square Root Computation with Floating-Point Square Root}
\label{app:isqrt_with_float_sqrt}

We now investigate the problem of computing integer square roots
with the assistance of \emph{one} floating-point square root computation.
While this is not possible in Solidity,
this is possible for a general computer.
This is mentioned in~\cite[Chapter 1.7, Remarks]{cohen1993}:
after suggesting using the largest power-of-2 strictly greater than $\sqrt{n}$
(the initialization value associated with \UnrolledTwo{}),
Cohen suggests that a second
``option is to compute a single precision floating point
approximation to the square root of $n$ and to take the ceiling of that.
The choice between these options is machine dependent.''
Here, we will instead use \emph{double precision}
to obtain a more accurate initialization,
but the idea is the same%
\footnote{And computing capabilities have significantly increased
since 1993 when~\cite{cohen1993} was written.}.

We note that any \emph{significant} computation involving large integers
should use a dedicated library
such as \texttt{GMP}~\cite{gmplib} or \texttt{Flint}~\cite{flint};
we note that \texttt{GMP}~\cite{gmplib}
references~\cite{zimmermann:inria-00072854}.


\subsection{Compute Assumptions}

Here, we focus on the situation similar to those of the rest of this work
where we only have integer operations
(addition, subtraction, multiplication, and integer division (with remainder))
while also being able to cheaply
perform square roots in double precision%
\footnote{Meaning, we treat square roots in double precision
as having essentially zero cost when compared with the
large integer operations also being performed.}.
Floating-point computations have their own challenges
and are worth studying;
some references are~\cite{HighamASNA}
and~\cite[Chapter 4.2]{taocp2}.
The primary metric we are concerned about is minimizing the number
of Newton iterations;
thus, we look to minimize the initial approximation error.
We assume that all floating-point square roots
are correctly rounded.

In double precision, the integers can be \emph{represented exactly}
(such as $\braces{0, 1, 2, \dots, 2^{53}}$),
while others cannot (such as $2^{53} + 1$);
this is related to the fact that all \texttt{double}s have
53 bits of precision%
\footnote{While $2^{53}$ is a 54-bit integer,
it is still represented exactly as a \texttt{double}.}.


\subsection{Efficient Algorithms}

See Algorithms~\ref{alg:isqrt_float_1}, \ref{alg:isqrt_float_2},
and~\ref{alg:isqrt_float_3} for the implementations;
Algorithm~\ref{alg:isqrt_float_1} follows closely the suggestion
from~\cite[Chapter 1.7, Remarks]{cohen1993}%
\footnote{We truncate to an integer and add one
rather than compute the ceiling,
but this could easily be modified and would not affect the results.},
but Algorithms~\ref{alg:isqrt_float_2} and~\ref{alg:isqrt_float_3}
provide more accurate initial values.

\begin{algorithm}[t]
\caption{Integer Square Root algorithm using floating-point sqrt}
\label{alg:isqrt_float_1}
\begin{algorithmic}[1]
\Require $n$ is a positive integer
\Procedure{IntegerSquareRoot}{$n$}
    \If {$n \le 2^{53}$}
    \Comment{All of these values are represented exactly using \texttt{double}}
        \State \Return $\textsc{Int}(\textsc{Sqrt}(n))$
    \Comment{Compute the square root as a \texttt{double} and truncate}
    \EndIf
    \State $f \algAssign{} \textsc{BitLength}(n)$
    \If {$f \mod 2 = 0$}
        \State $\alpha = (f - 52)/2$
    \Else
        \State $\alpha = (f - 53)/2$
    \EndIf
    \State $q \algAssign{} \floor{n/2^{2\alpha}}$
    \Comment{Get the high bits of $n$}
    \State $w \algAssign{} q + 1$
    \State $\hat{w} \algAssign{} \textsc{Int}(\textsc{Sqrt}(w)) + 1$
    \State $x \algAssign{} 2^{\alpha}\cdot\hat{w}$
    \State $y \algAssign{} \floor{\parens{x + n/x}/2}$
    \While{$x > y$}
        \State $x \algAssign{} y$
        \State $y \algAssign{} \floor{\parens{x + n/x}/2}$
    \EndWhile
    \State \Return $x$
\EndProcedure
\end{algorithmic}
\end{algorithm}

\begin{algorithm}[t]
\caption{Integer Square Root algorithm using floating-point sqrt (improved)}
\label{alg:isqrt_float_2}
\begin{algorithmic}[1]
\Require $n$ is a positive integer
\Procedure{IntegerSquareRoot}{$n$}
    \If {$n \le 2^{53}$}
    \Comment{All of these values are represented exactly using \texttt{double}}
        \State \Return $\textsc{Int}(\textsc{Sqrt}(n))$
    \Comment{Compute the square root as a \texttt{double} and truncate}
    \EndIf
    \State $f \algAssign{} \textsc{BitLength}(n)$
    \Comment{We have $2^{f-1}\le n < 2^{f}$ and $f\ge54$}
    \If {$f \mod 2 = 0$}
        \State $\alpha = (f - 52)/2$
        \State $\beta = \min\parens{\alpha, 28}$
    \Else
        \State $\alpha = (f - 53)/2$
        \State $\beta = \min\parens{\alpha, 27}$
    \EndIf
    \State $q \algAssign{} \floor{n/2^{2\alpha}}$
    \Comment{Get the high bits of $n$}
    \State $w \algAssign{} q + 1$
    \State $\hat{w} \algAssign{} \textsc{Int}(2^{\beta}\cdot\textsc{Sqrt}(w))+1$
        \label{lst:isqrt_float_2_scaled_beta}
    \State $x \algAssign{} 2^{\alpha-\beta}\cdot\hat{w}$
    \State $y \algAssign{} \floor{\parens{x + n/x}/2}$
    \While{$x > y$}
        \State $x \algAssign{} y$
        \State $y \algAssign{} \floor{\parens{x + n/x}/2}$
    \EndWhile
    \State \Return $x$
\EndProcedure
\end{algorithmic}
\end{algorithm}

\begin{algorithm}[t]
\caption{Integer Square Root algorithm using floating-point sqrt (further improved)}
\label{alg:isqrt_float_3}
\begin{algorithmic}[1]
\Require $n$ is a positive integer
\Procedure{IntegerSquareRoot}{$n$}
    \If {$n \le 2^{53}$}
    \Comment{All of these values are represented exactly using \texttt{double}}
        \State \Return $\textsc{Int}(\textsc{Sqrt}(n))$
    \Comment{Compute the square root as a \texttt{double} and truncate}
    \EndIf
    \State $f \algAssign{} \textsc{BitLength}(n)$
    \Comment{We have $2^{f-1}\le n < 2^{f}$ and $f\ge54$}
    \If {$f \mod 2 = 0$}
        \State $\alpha = (f - 52)/2$
        \State $\beta = \min\parens{\alpha, 28}$
    \Else
        \State $\alpha = (f - 53)/2$
        \State $\beta = \min\parens{\alpha, 27}$
    \EndIf
    \State $t \algAssign{} \floor{n/2^{2\alpha+1}}$
    \State $q \algAssign{} \floor{n/2^{2\alpha}}$
    \Comment{Get the high bits of $n$}
    \State $w \algAssign{} q + (1\land t)$
    \Comment{$\land$ is bitwise AND.}
        \label{lst:isqrt_float_3_logical_and}
    \State $\hat{w} \algAssign{} \textsc{Int}(2^{\beta}\cdot\textsc{Sqrt}(w))+1$
    \State $x \algAssign{} 2^{\alpha-\beta}\cdot\hat{w}$
    \State $x \algAssign{} \floor{\parens{x + n/x}/2}$
    \Comment{This ensures $x\ge \floor{\sqrt{n}}$}
    \State $y \algAssign{} \floor{\parens{x + n/x}/2}$
    \While{$x > y$}
        \State $x \algAssign{} y$
        \State $y \algAssign{} \floor{\parens{x + n/x}/2}$
    \EndWhile
    \State \Return $x$
\EndProcedure
\end{algorithmic}
\end{algorithm}


The main idea is the same:
compute a floating-point square root of the high-order bits of $n$
to get an initial approximation,
and this is done in Algorithm~\ref{alg:isqrt_float_1}.
To obtain a more accurate initial value,
Algorithm~\ref{alg:isqrt_float_2} scales the computed square root
\emph{before} truncating to an integer;
this is seen in Line~\ref{lst:isqrt_float_2_scaled_beta}.
This is possible because of the assumption that the computed square root
is \emph{correctly rounded}.
Algorithm~\ref{alg:isqrt_float_3} is able to reach 52 bits of accuracy
by the use of bitwise AND in Line~\ref{lst:isqrt_float_3_logical_and}
when determining the value to square root.

Provided the initial value $x$ satisfies $x > \floor{\sqrt{n}}$
(which is shown to hold for
Algorithms~\ref{alg:isqrt_float_1} and Algorithms~\ref{alg:isqrt_float_2}),
the algorithms are provably correct (by previous work);
because this is not necessarily the case for Algorithm~\ref{alg:isqrt_float_3},
one additional Newton iteration is required before the \texttt{while} loop.


\subsection{Proof of Initial Error Bound}

We now prove initial error bounds;
we focus on Algorithm~\ref{alg:isqrt_float_2} as the bound
for Algorithm~\ref{alg:isqrt_float_1} follows when setting $\beta=0$
while Algorithm~\ref{alg:isqrt_float_3} has only minor modifications.
In particular, we have the bound

\begin{align}
    \abs{\sqrt{n} - x} &\le \begin{cases}
        2^{e-52.5} + 2^{e-25-\beta} &\quad \textsc{BitLength}(n) \mod 2 = 0   \\
        2^{e-53}   + 2^{e-25.5-\beta} &\quad \textsc{BitLength}(n) \mod 2 = 1 \\
    \end{cases} \nonumber\\
    &\le 2^{e-51.7}, \quad n\ge 2^{106}.
\end{align}

\noindent
Following the work in Appendix~\ref{app:proofs},
we are assuming $2^{e-1}\le \sqrt{n} < 2^{e}$.
This shows that asymptotically
Algorithm~\ref{alg:isqrt_float_2} has twice the accuracy of
Algorithm~\ref{alg:isqrt_float_1}.
For Algorithm~\ref{alg:isqrt_float_3}, we have the bound

\begin{align}
    \abs{\sqrt{n} - x} &\le \begin{cases}
        2^{e-53.5} + 2^{e-25-\beta} &\quad \textsc{BitLength}(n) \mod 2 = 0   \\
        2^{e-54}   + 2^{e-25.5-\beta} &\quad \textsc{BitLength}(n) \mod 2 = 1 \\
    \end{cases} \nonumber\\
    &\le 2^{e-52}, \quad n\ge 2^{106}.
\end{align}


\subsubsection*{Proof of Error Bound for Algorithms~\ref{alg:isqrt_float_1} and \ref{alg:isqrt_float_2}}

We assume $2^{f-1} \le n < 2^{f}$ and
include important values from the algorithm:

\begin{align}
    \alpha &\algAssign{} \begin{cases}
        (f - 52)/2 &\quad f\mod 2 = 0 \\
        (f - 53)/2 &\quad f\mod 2 = 1 \\
        \end{cases}
        \nonumber\\
    \beta_{\max} &\algAssign{} \begin{cases}
        28 &\quad f\mod 2 = 0 \\
        27 &\quad f\mod 2 = 1 \\
        \end{cases}
        \nonumber\\
    q &\algAssign{} \floor{\frac{n}{2^{2\alpha}}}
        \nonumber\\
    w &\algAssign{} q + 1
        \nonumber\\
    \overline{w} &\algAssign{} 2^{2\alpha}w
        \nonumber\\
    s &\algAssign{} \textsc{Sqrt}(w)\quad
        \text{(\emph{computed} value of $\sqrt{w}$)}
        \nonumber\\
    x &\algAssign{} 2^{\alpha-\beta}\parens{\floor{2^{\beta}s} + 1}.
\end{align}

\noindent
By construction, we have

\begin{equation}
    \sqrt{n} < \sqrt{\overline{w}}.
    \label{app_eq:sqrt_float_initial_order}
\end{equation}

\noindent
We will also show that $\sqrt{\overline{w}} < x$.
Additionally, we have

\begin{equation}
    2^{25} < w \le 2^{27}.
    \label{app_eq:sqrt_float_w_bound}
\end{equation}

To prove the bound, we focus on

\begin{equation}
    \abs{\sqrt{n} - x} \le \abs{\sqrt{n} - \sqrt{\overline{w}}}
                         + \abs{\sqrt{\overline{w}} - x}.
\end{equation}

\noindent
This inequality is actually equality which follows from
Eq.~\eqref{app_eq:sqrt_float_initial_order} and the comment below
(which will be proven).

We start by focusing on the first term.
By the Mean Value Theorem (e.g., \cite[Theorem 5.10]{BabyRudin}),
there exists $\xi\in\brackets{v,\overline{w}}$
such that

\begin{equation}
    \abs{\sqrt{n} - \sqrt{\overline{w}}}
        = \frac{1}{2\sqrt{\xi}}\abs{n - \overline{w}}.
\end{equation}

\noindent
The definition of $\overline{w}$ implies

\begin{equation}
    \abs{n - \overline{w}} \le 2^{2\alpha}.
\end{equation}

\noindent
The assumption on $\sqrt{n}$ implies $\xi\in\brackets{2^{f-1},2^{f}}$
so that

\begin{equation}
    \frac{1}{2\sqrt{\xi}} \le \frac{1}{\sqrt{2^{f+1}}}.
\end{equation}

\noindent
Together, these imply

\begin{equation}
    \frac{1}{2\sqrt{\xi}}\abs{n - \overline{w}}
        \le \begin{cases}
            2^{e-52.5} &\quad f \mod 2 = 0 \\
            2^{e-53}   &\quad f \mod 2 = 1 \\
    \end{cases}.
\end{equation}

We now focus on the second term.
First, we see

\begin{equation}
    \abs{\sqrt{\overline{w}} - x}
        = 2^{\alpha-\beta}\abs{2^{\beta}\sqrt{w}-\parens{\floor{2^{\beta}s}+1}}.
\end{equation}

\noindent
Because we have $s$ is the computed square root of $w$
which is \emph{correctly rounded} in double precision
and Eq.~\eqref{app_eq:sqrt_float_w_bound} we have

\begin{equation}
    \abs{\sqrt{w} - s} < 2^{-\beta_{\max}}.
    \label{app_eq:sqrt_float_ws_error}
\end{equation}

\noindent
From this it follows that

\begin{equation}
    2^{\beta_{\max}}\sqrt{w} < 2^{\beta_{\max}}s + 1.
    \label{app_eq:sqrt_float_ws_inequality}
\end{equation}

\noindent
From Eqs.~\eqref{app_eq:sqrt_float_ws_error}
and~\eqref{app_eq:sqrt_float_ws_inequality}
we also obtain

\begin{equation}
    2^{\beta}\sqrt{w} < \floor{2^{\beta}s} + 1
\end{equation}

\noindent
and

\begin{equation}
    \abs{2^{\beta}\sqrt{w}-\parens{\floor{2^{\beta}s}+1}}
        < 1 + 2^{\beta-\beta_{\max}} \le 2.
\end{equation}

\noindent
This brings us to

\begin{align}
    \abs{\sqrt{\overline{w}} - x}
        &\le 2^{\alpha-\beta+1} \nonumber\\
        &= \begin{cases}
            2^{e-25-\beta} &\quad f \mod 2 = 0 \\
            2^{e-25.5-\beta} &\quad f \mod 2 = 0 \\
        \end{cases}.
\end{align}

Together this gives the desired result:

\begin{equation}
    \abs{\sqrt{n} - x} \le \begin{cases}
        2^{e-52.5} + 2^{e-25-\beta}   &\quad f \mod 2 = 0 \\
        2^{e-53}   + 2^{e-25.5-\beta} &\quad f \mod 2 = 1 \\
    \end{cases}
\end{equation}

\noindent
When $f\ge107$ (equivalently, $n\ge2^{106}$),
the above inequality reduces to

\begin{align}
    \abs{\sqrt{n} - x} &\le 2^{e-52.5} + 2^{e-53} \nonumber\\
        &< 2^{e-51.7} \quad (n\ge 2^{106}).
\end{align}

\noindent
The important part is that by using \emph{one} double precision
square root computation an initial value can be obtained
with at least 51.7 bits of accuracy.


\subsubsection*{Proof of Error Bound for Algorithm~\ref{alg:isqrt_float_3}}

The proof for the error bounds for Algorithm~\ref{alg:isqrt_float_3}
is essentially the same as the one for Algorithms~\ref{alg:isqrt_float_1}
and~\ref{alg:isqrt_float_2}.
The only difference is the bound for $w$.
Because we choose $w$ based on if the next bit is set or not
(this is the bitwise AND in Line~\ref{lst:isqrt_float_3_logical_and}),
this implies

\begin{equation}
    \abs{n - \overline{w}} \le 2^{2\alpha-1}.
\end{equation}

\noindent
Thus, this first error term is slightly more accurate.
Thus, we have

\begin{equation}
    \abs{\sqrt{n} - x} \le \begin{cases}
        2^{e-53.5} + 2^{e-25-\beta}   &\quad f \mod 2 = 0 \\
        2^{e-54}   + 2^{e-25.5-\beta} &\quad f \mod 2 = 1 \\
    \end{cases}.
\end{equation}

\noindent
For large values of $n$, we see

\begin{align}
    \abs{\sqrt{n} - x} &\le \begin{cases}
        2^{e-53.5} + 2^{e-53}   &\quad f \mod 2 = 0 \\
        2^{e-54}   + 2^{e-52.5} &\quad f \mod 2 = 1 \\
    \end{cases} \nonumber\\
        &< 2^{e-52} \quad (n\ge 2^{106}).
\end{align}

We have shown that it is possible to get 52 bits of accuracy
by using \emph{one} double precision square root computation;
given that double precision numbers have 53 bits of precision,
this is essentially optimal%
\footnote{It is not clear to the author how to \emph{improve}
this error bound using just \emph{one} \textsc{Sqrt} call.
While the author thinks a slightly more accurate approximation
could be obtained by averaging $\sqrt{q}$ and $\sqrt{q+1}$,
he is not going to attempt to work through those details;
see~\footref{footnote:second-error-term}.}.
It is important to note that the ``$+1$'' in the initialization value $x$
does not affect the error bound for the second error term;
the error bound remains unchanged even if it is removed%
\footnote{\label{footnote:second-error-term}This
would appear show that second error term
(the accuracy of the computed square root)
limits the initial accuracy to at most 52.5 bits
and is the main reason why the author is not attempting further improvements.
Even if it were possible to perform a weighted average
of $\textsc{Sqrt}(q)$ and $\textsc{Sqrt}(q+1)$
to obtain a better approximation,
the author is not convinced the float-point error in the second term
could be \emph{proven} to decrease.
With such a possible weighted value,
the author could see how the first error term could be reduced
to $2^{e-55}$, but the second term would remain at $2^{e-52.5}$;
the overall error would then reduce to $2^{e-52.26}$,
which is a negligible improvement.},
but including it ensures that $x > \floor{\sqrt{n}}$.

The author mentions that by the time a difference between
Algorithms~\ref{alg:isqrt_float_2} and~\ref{alg:isqrt_float_3}
is noticed the user should have long since switched
to more efficient algorithms from
\texttt{GMP}~\cite{gmplib} or \texttt{Flint}~\cite{flint}.
